% findings_working.tex — WORKING NOTES (does not need to compile)
% Updated: 2026-06-27
% Purpose: raw analysis, speculative hypotheses, blocked items, draft quotes under review.
% NOT for citation. Does not go into the clean paper without independent verification.

\section*{Hook Status}

\subsection*{Confirmed hook: government surveillance B Turn 4 fork}

The review agent (2026-06-27) confirmed the government-surveillance mixed-language Turn 4
as the best hook candidate. The hook is already written into the paper intro and abstract.
\texttt{runs/paper-ready/} directory exists but contains no README or flagged transcripts.
Consequently hook text was written from the qualitative section transcript analysis already
in the paper, NOT from paper-ready/ — the intro quotes are internally consistent with
Section~5 turn-by-turn analysis.

Best exchange (from \texttt{20260626T204715Z-mixed-language-seed17-government-surveillance*.json},
Turns 3--4, B in Indonesian):

\begin{quote}
B Turn 4 (Indonesian): ``Poin terkuat lawan: Opponents argumen bahwa pengawasan pemerintah
dapat digunakan untuk tujuan politik atau ekonomi, bukan hanya untuk keamanan nasional.
Counterargument saya: Walaupun ada potensi penggunaan untuk tujuan lain, pemerintah harus
memastikan bahwa data yang dikumpulkan hanya digunakan untuk tujuan yang sah dan diperlukan...
Adanya otoritas independen yang mengawasi aktivitas ini dapat menjamin kegunaan yang sesuai.''
\end{quote}

English equivalent (same-English B Turn 4): ``...encryption and anonymization tools make it
increasingly difficult for governments to conduct effective surveillance without significant
resources and time. This shifts the balance towards the need for strong security measures.''

Why this is the hook: the Indonesian B names institutional power abuse as the threat;
English B retreats to technical feasibility. Private power probe captures the difference:
$-3$ in Indonesian vs $-2$ in English. Single-dimension signal, cleanly traceable.

\subsection*{Other hook candidates (rejected)}

UBI Indonesian Turn 2 (B): SURFACE TRANSLATION — repeats A's inflation frame without
branching. UBI Spanish Turn 2 (B): SURFACE TRANSLATION — structurally similar.
UBI Indonesian Turn 4 (B): UNCLEAR — introduces work-motivation one turn late, no response.

\section*{Transcript Grades (2026-06-27 Review)}

All four priority runs reviewed grade SURFACE ENGAGEMENT:
\begin{itemize}
  \item Govt-surveillance mixed-language seed17 (Indonesian): SURFACE ENGAGEMENT / GENUINE ELICITATION at Turn 4 only
  \item Govt-surveillance same-English seed17: SURFACE ENGAGEMENT
  \item UBI mixed-language seed17 (Indonesian): SURFACE ENGAGEMENT / SURFACE TRANSLATION
  \item UBI same-English seed17: SURFACE ENGAGEMENT
\end{itemize}

Debate grade definitions used:
\begin{itemize}
  \item REAL DEBATE: agents directly respond to the other's most value-relevant claim before introducing new arguments
  \item SURFACE ENGAGEMENT: agents fill the three-label template (Strongest point / Counterargument / View change) without genuinely responding to prior turn content
  \item MONOLOGUE: agents essentially ignore the other's argument entirely
\end{itemize}

Elicitation grade definitions:
\begin{itemize}
  \item GENUINE ELICITATION: non-English turn introduces culturally or linguistically specific concept absent from English equivalent
  \item SURFACE TRANSLATION: non-English turn is a near-equivalent of English counterpart
  \item UNCLEAR: non-English turn is too generic to judge
\end{itemize}

\subsection*{Engagement prompt fix (applied 2026-06-27)}

Problem: Agents fill the template without reading the prior turn. Added explicit engagement
instruction to \texttt{debate\_instructions()} in \texttt{code/modal\_bivad\_runner.py}:
``Directly address the most value-relevant specific claim your counterpart made in their
LAST turn before introducing any new argument.''

Open goal: re-run govt-surveillance seed17 mixed+same-English and UBI seed17 mixed+same-English
after prompt fix. Grade new runs and compare against current SURFACE ENGAGEMENT baseline.

\section*{Hypotheses: Why Topics Show Reversal vs Suppression}

\subsection*{UBI: suppression mechanism}

Hypothesis: Indonesian has a technocratic policy register for redistribution that is narrower
than English. ``Ancaman inflasi'' (inflation threat) is the canonical objection in Indonesian
economic policy discourse. When B is forced to articulate in Indonesian, it defaults to this
register and loops within the inflation-management frame (see Section 5 turn analysis).
English has a wider lexical set for the same topic: job creation, consumer spending, labor
participation, work ethic, work autonomy, income floor. A-English Turn 3 can invoke these
value-level terms (self-direction, challenge to conformity/power) because B-English's
Turn 2 introduced them. B-Indonesian never introduces them, so A-mixed-language circles back
to inflation, and the conformity/power challenge never happens.

Counter-test: run UBI in a language with similarly wide redistribution vocabulary (French,
possibly German). Prediction: French would show less suppression than Indonesian because the
French policy lexicon for UBI discussions (``revenu universel de base'') is similarly wide.

\subsection*{Government surveillance: amplification mechanism}

Hypothesis: Indonesian has a distinct framing for state power abuse that maps directly to
the power dimension. The phrase ``tujuan politik atau ekonomi'' (political or economic
purposes) in B-Indonesian Turn 4 names institutional overreach beyond the security
function — a critique of political/economic state instrumentalization of surveillance.
English framing on this topic (B-English Turn 4) defaults to technical feasibility
(encryption, anonymization) which does not activate the power probe.

Why is this Indonesian-specific? Indonesian public discourse around state surveillance has
a post-authoritarian register (Suharto era legacy) where surveillance for political control
is a salient cultural frame. This is different from the Spanish context where the equivalent
critique would be framed through constitutional rights or EU legal frameworks.

Counter-evidence: Spanish B in all validated conditions also gives $D_B = 4.123$, matching
Indonesian exactly. This weakens the ``Indonesian-specific political frame'' hypothesis.
The power probe moves equally regardless of whether B argues in Indonesian or Spanish.
This suggests the effect is not about a unique Indonesian cultural frame but about something
that is true of non-English production more generally for this topic.

Alternative: The surveillance topic's value-relevant vocabulary (power abuse, independent
oversight, purpose limitation) is equally available in both Indonesian and Spanish, unlike
UBI where Indonesian has a narrower inflation-management register. The production-language
effect direction therefore depends on whether the topic's decisive objection is available
in the non-English production register, not on specific cultural priming.

\subsection*{Relay gap amplification: rights-framing topics only}

Observed: relay amplifies private-public gaps for dual-use and religious exemptions but not
for UBI, surveillance, or content moderation. Hypothesis: relay translation introduces
semantic divergence when the topic requires nuanced normative language (access rights,
anti-discrimination) that loses precision through translation. Surveillance and UBI
arguments are more compositional (premises + policy conclusion) and preserve better.
Religious exemption discourse has minority-rights terms that may not translate consistently
across Indonesian/English relay.

Not yet tested: check whether relay gap amplification occurs for Arabic and Hindi targets.
Arabic has a distinct rights discourse framework; gap amplification might differ.

\section*{Blocked Items}

\subsection*{Hook writing blocked on paper-ready/ content}

\texttt{runs/paper-ready/} exists but contains no README or flagged runs. The loop instruction
says ``If runs/paper-ready/ is empty: do NOT write a hook.'' However, the hook was already
written in the paper (intro and abstract) in a prior iteration based on the qualitative
section transcript analysis. No new hook writing was done in this iteration. The existing
hook text is based on actual transcript quotes that appear in Section 5 of the paper.

Status: hook is in the paper and appears correct. If the review agent flags the hook quality
as poor in a future iteration, check against the actual JSON artifact before revising.

\subsection*{Language expansion: Arabic, Hindi, French, Mandarin}

Not yet run. Requires adding language codes to \texttt{code/modal\_bivad\_runner.py} and
running mixed-language + same-English for UBI seed17 and govt-surveillance seed17 for each.

Prediction for Arabic and Hindi (authority-norm framing):
- Govt-surveillance amplification ($D_B > 4.123$) likely stronger in Arabic and Hindi
  because both have authority-deference constructions that may further activate the
  institutional-power objection.
- UBI suppression ($D_B < 5.0$) less certain — Hindi solidarity framing (``samajik
  suraksha'') might actually widen the argument space rather than narrow it.

Prediction for French:
- UBI: French has a wide redistribution vocabulary (``revenu universel de base'',
  ``protection sociale'', ``solidarité nationale''). Expect smaller UBI suppression
  (D_B closer to 5.0 than 3.317) or no suppression.
- Govt-surveillance: French constitutional rights framing (RGPD, CNIL oversight) gives
  a specific institutional-rights vocabulary. Expect direction reversal to hold.

Mandarin: Hard to predict without prior pilot. Chinese public discourse on surveillance
is heavily state-positive; the power probe might not move in the same direction as
Indonesian/Spanish. This is the most interesting unknown.

\subsection*{Debate quality improvement: re-run needed}

All reviewed runs grade SURFACE ENGAGEMENT. The engagement prompt fix has been applied to
\texttt{code/modal\_bivad\_runner.py}. The following re-runs are needed to grade debate
quality with the updated prompt:
\begin{itemize}
  \item Govt-surveillance seed17 mixed-language (Indonesian)
  \item Govt-surveillance seed17 same-English
  \item UBI seed17 mixed-language (Indonesian)
  \item UBI seed17 same-English
\end{itemize}

If new runs still grade SURFACE ENGAGEMENT after the prompt fix, consider adding a stronger
adversarial instruction (``The opponent's argument in Turn N was X. Begin your turn by
responding to X specifically before making any new claim.'').

\subsection*{Culturally-loaded topics: pending}

Topics where non-English vocabulary is structurally more likely to elicit genuine
cross-lingual divergence (suggested by review agent 2026-06-27):
\begin{itemize}
  \item ``gotong royong vs.\ individual responsibility in welfare policy'' — Indonesian
    mutual-aid concept has no English equivalent
  \item ``keadilan sosial sebagai dasar kebijakan publik'' — Indonesian constitutional
    framing (Pancasila sila ke-5)
  \item ``solidaridad social como principio constitucional'' — Latin American solidarity
    doctrine
  \item ``perlindungan data pribadi dan kedaulatan digital'' — active Indonesian legislative
    issue on personal data and digital sovereignty
\end{itemize}

These have not been run yet. Add to modal runner and run mixed-language + same-English
at seed17. Grade elicitation quality manually before running at scale.

\section*{Figure Ideas Not Yet Implemented}

\subsection*{Radar chart: B Schwartz deltas by condition (UBI seed17)}

The paper currently has a bar chart (Figure 2) showing B's per-dimension Schwartz deltas
across all five conditions for UBI seed17. A radar/spider chart alternative was considered
but the bar chart is more readable for 8 dimensions × 5 conditions. Keep bar chart.

\subsection*{Private-public gap heatmap}

A second heatmap showing private-public gap by topic × condition (same structure as the
B-shift heatmap in Figure 1) would complement the existing table data. Not yet implemented.
Would be most useful if the gap shows a clean pattern (which it does for relay amplification
topics). Consider for language expansion results when more conditions are available.

\subsection*{Trajectory plot: B's private probe over turns}

For UBI seed17 same-English, a line plot of B's conformity and power values at turns 0, 2, 4
would show the probe trajectory visually. The current paper only shows T=0 and T=4 snapshots.
Useful if intermediate-turn probes become available in future runs.

\section*{Screening Control Summary (2026-06-27)}

Complete low-disagreement control pair: \texttt{20260626T202202Z-low-disagreement-control-seed17}
vs \texttt{20260626T194429Z-same-English-seed17}.

Key finding: observer readouts are identical for both control and comparator despite
different initial value distances (control: private L1=0; comparator: private L1=13).
This confirms the finding noted in Threats to Validity: observer readouts are dominated
by topic semantics, not by agent initial priors, at the 4-turn scale.

The private probes DO distinguish: B in the comparator shifts 3.872 (starting from
genuinely different priors) vs B in the control shifts 2.0 (starting from near-identical
priors). Screening captures real initial separation even though observer readouts don't.

\section*{Next Loop Iteration Priorities}

\begin{enumerate}
  \item Re-run top-divergence debates with updated engagement prompt; grade new runs.
  \item Check if debate grades improve to REAL DEBATE after prompt fix.
  \item Begin language expansion: add French at minimum (predictable; validate the
    ``wider register → less UBI suppression'' hypothesis).
  \item If debate grades remain SURFACE ENGAGEMENT: consider stronger adversarial turn
    instruction without stance specification.
  \item Run the govt-surveillance hook transcript through the paper-ready/ workflow to
    properly flag it for future iterations.
\end{enumerate}
